\documentclass[12pt]{article}
\usepackage{amsmath,amssymb}
\usepackage[margin=1in]{geometry}
\usepackage{setspace}
\onehalfspacing

\title{Origins of Modern Economics and Econometrics in  the Mechanics of Lagrange, Hamilton, and Laplace: A Sad Tale about Forgetting and Ignoring}

\author{Thomas J. Sargent}
\date{}

\begin{document}

\maketitle

\section{Introduction}

Three representations of classical mechanics  influenced the mid-twentieth century intellectual giants who created modern economic theory and econometrics: \begin{itemize}
\item Lagrange's variational formulation
\item Hamilton's canonical and symplectic formulation
\item Laplace's deterministic differential-equation formulation
\end{itemize}


\noindent These  three %technical machineries
 look at the same object from  different perspectives.  
 %Mean field games inherit structural elements from all three:
\begin{itemize}
\item  Lagrange, the variational structure and optimality,
\item  Hamilton, a set of  canonical equations and Hamilton--Jacobi theory,
\item  Laplace, equations that describe deterministic evolution of distributions.
\end{itemize}
Their %Lagrange, Hamilton, and Laplace 
 complementary formulations of classical mechanics   shaped modern economics.
% from optimal control to macroeconomic dynamics and strategic interaction in large populations.


This document presents each formulation with its key equations, derives the Euler--Lagrange equations, introduces canonical transformations, 
%explains Pontryagin's Maximum Principle,
 and notes how mean field games uses ideas from classical mechanics.


%\section{Part I: Three Formulations of Classical Mechanics}

\subsection*{Lagrange's Analytic Mechanics}

Let $q = (q_1, \ldots, q_n)$ be generalized coordinates and
\[
L(q, \dot{q}, t) = T(q, \dot{q}) - V(q)
\]
the Lagrangian.
The action functional is
\[
S[q] = \int_{t_0}^{t_1} L(q(t), \dot{q}(t), t) \, dt.
\]

\subsubsection*{Euler--Lagrange Equations}

Consider a variation $q_i(t) + \varepsilon \eta_i(t)$ with $\eta_i(t_0) = \eta_i(t_1) = 0$. Then
\[
\delta S = \varepsilon \int_{t_0}^{t_1} \left( \frac{\partial L}{\partial q_i} \eta_i + \frac{\partial L}{\partial \dot{q}_i} \dot{\eta}_i \right) dt.
\]
Integrate the second term by parts:
\[
\int_{t_0}^{t_1} \frac{\partial L}{\partial \dot{q}_i} \dot{\eta}_i \, dt = \left[ \frac{\partial L}{\partial \dot{q}_i} \eta_i \right]_{t_0}^{t_1} - \int_{t_0}^{t_1} \frac{d}{dt} \left( \frac{\partial L}{\partial \dot{q}_i} \right) \eta_i \, dt.
\]
The boundary term vanishes, so
\[
\delta S = \varepsilon \int_{t_0}^{t_1} \left( \frac{\partial L}{\partial q_i} - \frac{d}{dt} \left( \frac{\partial L}{\partial \dot{q}_i} \right) \right) \eta_i(t) \, dt.
\]
Since this must hold for all $\eta_i$, we obtain the \textbf{Euler--Lagrange equations}:
\[
\frac{d}{dt} \left( \frac{\partial L}{\partial \dot{q}_i} \right) - \frac{\partial L}{\partial q_i} = 0.
\]
%
%\subsubsection*{Applications in Economics}
%
%Lagrangian methods appear in: dynamic programming (Bellman equations), optimal growth theory (Ramsey--Cass--Koopmans), continuous-time consumption--investment problems, and mechanism design with dynamic constraints.
%Euler--Lagrange equations appear as  first-order conditions for optimal control problems.


\subsection*{Hamilton's Canonical Mechanics}

Define conjugate momenta:
\[
p_i = \frac{\partial L}{\partial \dot{q}_i}.
\]
The Hamiltonian is the Legendre transform of the Lagrangian:
\[
H(q, p, t) = \sum_{i=1}^n p_i \dot{q}_i - L(q, \dot{q}, t).
\]
Hamilton's equations are:
\[
\dot{q}_i = \frac{\partial H}{\partial p_i}, \quad \dot{p}_i = -\frac{\partial H}{\partial q_i}.
\]

\subsubsection*{Canonical Transformations}

A transformation $(q, p) \mapsto (Q, P)$ is \emph{canonical} if it preserves the \emph{symplectic} form:%
%A transformation $(q, p) \mapsto (Q, P)$ is \emph{canonical} if it preserves the \emph{symplectic} form:
\footnote{The wedge symbol $\wedge$ denotes the wedge product (or exterior product) from differential geometry. The expression $dq_i \wedge dp_i$ is an antisymmetric bilinear form satisfying $dq_i \wedge dp_i = -dp_i \wedge dq_i$ and $dq_i \wedge dq_i = 0$. When evaluated on two tangent vectors, it measures the signed area of the parallelogram they span in the $(q_i, p_i)$ plane. The condition $\sum_i dq_i \wedge dp_i = \sum_i dQ_i \wedge dP_i$ thus says that canonical transformations preserve the oriented area structure of phase space, ensuring that Hamilton's equations maintain their form in the new coordinates.}
%\[
%\sum_i dq_i \wedge dp_i = \sum_i dQ_i \wedge dP_i.
%\]
%
\[
\sum_i dq_i \wedge dp_i = \sum_i dQ_i \wedge dP_i.
\]
Such a transformation preserves Hamilton's equations. Canonical transformations are generated by functions $F$ that satisfy
\[
p_i = \frac{\partial F}{\partial q_i}, \quad Q_i = \frac{\partial F}{\partial P_i}.
\]
%
%\subsubsection*{Applications in Economics}
%
%Canonical transformations underlie: dynamic programming and costate transformations, asset pricing models with Hamilton--Jacobi--Bellman equations, symplectic integrators  in macroeconomic simulations, and transformations that simplify optimal control problems.
%They are  backbones of continuous-time dynamic optimization.


\subsection*{Laplace's Deterministic Mechanics}

Laplace's formulation is based on deterministic ODEs. For $N$ bodies define:
\[
m_i \ddot{x}_i = \sum_{j \neq i} G \frac{m_i m_j (x_j - x_i)}{\|x_j - x_i\|^3}.
\]
For Laplace, the current  state vector of a mechanical system determines its  future and its past. 
%Explanation is deductive: laws + initial conditions. No metaphysical ``cause'' is needed.
%
%\subsubsection*{Remarks for Economics}
%
%Laplace's determinism parallels: rational expectations equilibria, structural econometric models, and deterministic laws of motion for distributions (e.g., Kolmogorov forward equations).
%
%
%\subsection{Pontryagin's Maximum Principle}
%
%Consider a control system:
%\[
%\dot{x}(t) = f(x(t), u(t)), \quad x(0) = x_0,
%\]
%with objective
%\[
%J(u) = \int_0^T L(x(t), u(t)) \, dt + \Phi(x(T)).
%\]
%
%Define the Hamiltonian:
%\[
%H(x, u, \lambda) = L(x, u) + \lambda^\top f(x, u).
%\]
%
%Pontryagin's Maximum Principle states: A necessary condition for optimality is that there exists a costate $\lambda(t)$ such that
%\[
%\dot{x}(t) = \frac{\partial H}{\partial \lambda}, \quad \dot{\lambda}(t) = -\frac{\partial H}{\partial x},
%\]
%and the optimal control satisfies
%\[
%u(t) \in \arg\max_v H(x(t), v, \lambda(t)).
%\]
%
%\subsubsection{Remarks for Economics}
%
%PMP is the foundation of: optimal growth theory, continuous-time macroeconomics, dynamic contract theory, dynamic mechanism design, optimal taxation (Mirrlees, Chamley--Judd), and stochastic control in finance.
%
%It is the economic analogue of Hamiltonian mechanics.
%



\section*{Exorcising a Ghost named Cause}

\subsection*{Kepler versus Newton: Description versus Structure}

The distinction between descriptive regularities and structural laws has deep roots in the history of science, and it was explicitly invoked by pioneers of modern economics and game theory to clarify what economic theory should aspire to achieve.

\subsubsection{von Neumann and Morgenstern on Scientific Method}

In the opening chapter of their 1944 \emph{Theory of Games and Economic Behavior}, John von Neumann and Oskar Morgenstern draw a sharp distinction between two approaches to scientific explanation \cite{vonNeumannMorgenstern1944}:

\textbf{The Keplerian approach:} Johannes Kepler observed the motions of planets and formulated mathematical laws describing their orbits---ellipses with the sun at one focus, areas swept out in equal times, and a relationship between period and distance. These were \emph{descriptive regularities}, empirically accurate but lacking a deeper theoretical foundation. Kepler's laws told us \emph{what} happens, but not \emph{why}.

\textbf{The Newtonian approach:} Isaac Newton provided a deeper explanation by deriving Kepler's laws from a small set of fundamental principles: the law of universal gravitation and the laws of motion. Newton's theory was \emph{structural} in the sense that it identified underlying mechanisms (gravitational force, inertia) from which the observed regularities could be deduced. Moreover, Newton's laws allowed predictions in novel situations where Kepler's laws did not apply---such as the motion of comets, the precession of the equinoxes, and the tides.

Von Neumann and Morgenstern argued that economics in the early twentieth century was still largely in its Keplerian stage: economists had identified various empirical regularities (supply and demand curves, business cycles, marginal propensities to consume) but lacked a unified theoretical framework from which these could be derived. Their project in game theory was to provide a Newtonian foundation for economics---a set of principles about rational behavior under strategic interaction from which economic phenomena could be systematically explained.

\subsubsection{Koopmans on ``Measurement Without Theory''}

Three years later, Tjalling C.\ Koopmans made a related argument in his famous 1947 review article ``Measurement Without Theory'' \cite{Koopmans1947}, which critiqued the empirical methodology of the National Bureau of Economic Research. Koopmans argued that purely descriptive statistical work---measuring correlations, identifying cycles, constructing indices---was insufficient for scientific progress in economics.

Koopmans insisted that measurement must be guided by theory. Without a structural model that specifies the behavioral relationships and constraints governing economic agents, empirical regularities are fragile and unreliable. Correlations observed in one historical period may vanish or reverse when policies or institutions change, precisely because the descriptive relationships are not grounded in invariant structural primitives.

The Koopmans critique anticipated what would later be formalized as the Lucas critique: reduced-form relationships are not stable under policy interventions because they confound structural parameters (preferences, technologies) with equilibrium outcomes (prices, allocations) that depend on the policy regime.

\subsubsection{The Kepler--Newton Analogy in Econometrics}

The Kepler--Newton distinction maps directly onto the contrast between reduced-form and structural-form econometrics:
\begin{itemize}
\item \textbf{Kepler $\approx$ Reduced form:} Empirical regularities that describe relationships among observables without identifying the underlying mechanisms. Examples: Phillips curves, consumption functions estimated by OLS, VAR models, diff-in-diff estimates of treatment effects.

\item \textbf{Newton $\approx$ Structural form:} Models that specify primitives (preferences, technologies, information, constraints) and derive equilibrium relationships. Examples: DSGE models, structural estimation of dynamic discrete choice models, calibrated general equilibrium models.
\end{itemize}

The advantage of the Newtonian/structural approach is \emph{invariance}. Just as Newton's laws apply to planets, comets, and terrestrial projectiles alike, structural economic models should yield reliable predictions under policy interventions or environmental changes not present in the historical data. The reduced-form approach, like Kepler's laws, may describe the data well within the sample but fails when the underlying structure changes.

\subsubsection{Cowles Commission and the Structural Program}

The Cowles Commission, under the leadership of Marschak and Koopmans, took up this challenge. Their goal was to develop econometric methods that could identify structural parameters from observational data, using a priori theoretical restrictions to achieve identification. This project was self-consciously modeled on the transition from Kepler to Newton in physics.

The Cowles Commission understood that:
\begin{enumerate}
\item Economic relationships are not fixed laws of nature but equilibrium conditions among optimizing agents.
\item To predict the effects of policy interventions, one must identify the invariant primitives (utility, technology) that govern behavior across different regimes.
\item Identification requires strong theoretical assumptions---what Hurwicz called ``a priori information''---because data alone cannot distinguish structural from reduced-form parameters.
\end{enumerate}

From this perspective, the modern causal inference literature represents a retreat to the Keplerian stage. By focusing on descriptive treatment effects estimated from natural experiments, without embedding these effects in a structural model, the literature sacrifices the generality and predictive power that comes from understanding the underlying mechanisms.

This is not to say that descriptive empirical work has no value---Kepler's laws were an essential step toward Newton's theory. But as von Neumann, Morgenstern, and Koopmans recognized, description is not the ultimate goal of science. The goal is structural understanding that supports reliable extrapolation and counterfactual reasoning.


\subsubsection*{Laplace's Vision}

Laplace's reformulation of classical mechanics is historically significant because it removes the notion of \emph{cause} from the explanatory structure of physics. In the \emph{M\'ecanique c\'eleste}, Newtonian mechanics is recast as a system of differential equations governing the evolution of celestial bodies. The state of the universe at a given instant---positions and velocities of all bodies---is sufficient to determine both its future and its past.

This perspective underlies the famous thought experiment now known as Laplace's demon: a hypothetical intellect that, knowing all forces and positions at one moment, could compute the entire temporal evolution of the universe. In this framework, explanation is not causal in a metaphysical sense; it is deductive. The present state, together with the laws of motion, \emph{is} the explanation.

Laplace avoids causal language. The laws of motion are expressed as relations among variables and their derivatives. Once initial conditions are specified, the system evolves uniquely. According to Laplace, if you want to say "cause" you should remember that it makes no sense to say that this variable (a particular component of the state? Some function of the state vector?) ``causes'' or is an ``effect of'' another variable:

\begin{quote}
``Nous devons envisager l'\'etat pr\'esent de l'univers comme l'effet de son \'etat ant\'erieur, et comme la cause de celui qui va suivre.''

\medskip
\raggedleft --- Laplace, \emph{Essai philosophique sur les probabilit\'es}
\end{quote}

\begin{quote}
``We must consider the present state of the universe as the effect of its prior state, and as the cause of the one which is to follow.''
\medskip

 \raggedleft ---Laplace, A Philosophical Essay on Probabilities 

\end{quote}


\subsubsection*{Redundancy of ``Causes''}
%
%\textbf{Determinism replaces causal explanation.} Laplace's mechanics is predictive rather than causal. The explanatory burden is carried by the differential equations themselves. The present state determines the future and the past; causal talk becomes redundant.

Saying that  ``this causes that'' is just  loose talk.  In Laplace's formulation, the present state tells everything pertinent about the past and about the future. The laws of motion tie together variables and their rates of change. It is at best redundant and usually  misleading  to say more.
Thus, when Napoleon Bonaparte asked im  about the role that God played in his celestial mechanics, Laplace 
told Napoleon  that he ``had no need of that hypothesis.'' That was Laplace's polite way of saying -- ``Just look at the equations and you can figure that out for yourself.''\footnote{With his answer to Napoleon's question, Laplace meant to sweep away over two thousand years of philosophers' having taken Aristotle's proof of the existence of God as a necessary "first cause''.}

Bertrand Russell summarized the situation in the following way.

\begin{quote}
There are some respects in which the concepts of modern theoretical physics differ from those of the Newtonian system. To begin with, the conception of `force', which is prominent in the seventeenth century, has been found to be superfluous. `Force', in Newton, is the cause of change of motion, whether in magnitude or direction. The notion of cause is regarded as important, and force is conceived imaginatively as the sort of thing that we experience when we push or pull.
%
%\medskip
%\raggedleft --- Bertrand Russell, \emph{History of Western Philosophy}
%\end{quote}
%
%\begin{quote}
Gradually it was found that all the equations could be written down without bringing in forces. What was observable was a certain relation between acceleration and configuration; to say that this relation was brought about by the intermediacy of `force' was to add nothing to our knowledge. Observation shows that planets have at all times an acceleration towards the sun, which varies inversely as the square of their distance from it. \ldots\ The modern physicist, therefore, merely states formulae which determine accelerations, and avoids the word `force' altogether. `Force' was the faint ghost of the vitalist view as to the causes of motions, and gradually the ghost has been exorcized.

\medskip
\raggedleft --- Bertrand Russell, \emph{History of Western Philosophy}
\end{quote}


%\subsection*{From Causal ``Agent'' to Function of the State Vector}

In Laplace's representation of classical mechanics, Newton's concept of ``force'' is demoted from a causal agent to a mathematical function appearing on the right-hand side of a differential equation.
%
%\subsection*{The Translation from Newton to Laplace}
%
%\textbf{Newton's formulation:}
In Newton's formulation\footnote{\textbf{Tom: correct the history of thought -- say where equation actually first appeared.}}
\[
\mathbf{F} = m\mathbf{a}
\]
Here, force is conceived as something that ``produces'' or ``causes'' acceleration. Force has a metaphysical status as the agent of change.

By way of contrast, in Laplace's formulation the state of a mechanical system is specified by positions and velocities: $\mathbf{x}(t)$ and $\mathbf{v}(t) = \dot{\mathbf{x}}(t)$. The system evolves according to:
\[
\frac{d\mathbf{x}}{dt} = \mathbf{v}
\]
\[
\frac{d\mathbf{v}}{dt} = \mathbf{f}(\mathbf{x}, \mathbf{v}, t)
\]
where $\mathbf{f}(\mathbf{x}, \mathbf{v}, t) = \mathbf{F}(\mathbf{x}, \mathbf{v}, t)/m$ is simply the acceleration as a function of the current state.

Thuss, in Laplace's framework, Newton's ``force'' is just the vector-valued function $\mathbf{f}(\mathbf{x}, \mathbf{v}, t)$ that specifies how velocity changes with time. It is a component of the differential equation, not a separate physical entity.
Consider a few examples, 
%
%\textbf{Examples:}

\begin{enumerate}
\item \textbf{Gravitational force near Earth's surface:}
\begin{itemize}
\item Newton: $F = -mg$
\item Laplace: $\displaystyle\frac{d^2x}{dt^2} = -g$
\end{itemize}

\item \textbf{Harmonic oscillator (spring force):}
\begin{itemize}
\item Newton: $F = -kx$
\item Laplace: $\displaystyle\frac{d^2x}{dt^2} = -\omega^2 x$ where $\omega^2 = k/m$
\end{itemize}

\item \textbf{Planetary motion (inverse square law):}
\begin{itemize}
\item Newton: $\mathbf{F} = -\dfrac{GMm}{r^2}\hat{\mathbf{r}}$
\item Laplace: $\displaystyle\frac{d^2\mathbf{r}}{dt^2} = -\frac{GM}{r^2}\hat{\mathbf{r}}$
\end{itemize}
\end{enumerate}

%\subsection*{The Philosophical Shift}

In Laplace's formulation:
Force is not an independent input causing motion---it is a function of the state vector.
There is just a mathematical rule describing how the state at time $t$ shapes future values of the state, and how it was shaped by past values of the state.
%, here is how it changes.
Thus, 
knowing the state and the function $\mathbf{f}$ lets us compute all past and future states.
We can note that ``acceleration depends on position and velocity as described by this function'' without ever mentioning force.


This is of course  what Bertrand Russell meant when he said that force is ``the faint ghost of the vitalist view'' that has been ``exorcised'' from modern physics. In Laplace's deterministic picture, forces are merely convenient names for particular functions appearing in differential equations, not independent metaphysical causes that make things happen.\footnote{For a physicist, ``metaphysics'' is a term of contempt.}

%
%\subsection{Laplace's Role in Exorcising ``Cause''}
%
%Seen against this background, Laplace completes a trajectory that runs from Newton's causal forces, through d'Alembert's equilibrium principle and Lagrange's analytic elimination of forces, to Hamilton's geometric evolution. Laplace's contribution is to insist that the present state and the laws of motion are sufficient for explanation. The concept of cause is not denied, but it is rendered superfluous within the formalism.
%

\section{Part III: Implications for Mathematical Economics}

\subsection{Laplace, Structure, and Economic Modeling}

Laplace's reformulation of mechanics influenced the 20th century giants who invented mathematical economics and post R.A.\ Fisher single-equation fertilizer frequentist tests. Thus, Marschak, Koopmans, and the Cowles Commission, cast their theories in terms of systems of stochastic difference or differential equations implied by solutions of a coherent set of optimization problems being solved by a precisely specified collection of decision makers, and explanation consists in showing how observed variables satisfy these structural relations under rational behavior and equilibrium conditions.

In this sense, the Laplacian move away from metaphysical cause toward structural, law-governed evolution anticipates the modern economist's reliance on structural equations, state variables, and dynamic laws.


\subsection{Mean Field Games and Laplacian Equilibrium}

Recent work on mean field games (MFGs) provides a striking modern analogue of Laplace's vision. In a continuous-time MFG, equilibrium is defined by a pair of coupled partial differential equations:
\begin{itemize}
\item a Hamilton--Jacobi--Bellman (HJB) equation for the value function of a representative infinitesimal agent, and
\item a Fokker--Planck (FP) equation describing the evolution of the distribution of agents.
\end{itemize}

Given the initial distribution and the structural primitives of the model, the entire path of the system is determined by these equations. The equilibrium is a fixed point: agents optimize given the distribution, and the distribution evolves according to the optimal controls.

This structure is deeply Laplacian. The distribution of agents at time $t$ plays the role of the ``state of the universe,'' and the coupled HJB--FP system provides the deterministic laws of motion. No causal narrative is required; the explanation lies in the structural equations and initial conditions.

From the perspective of Marschak and Koopmans, mean field games exemplify the shift from causal stories to structural, law-governed dynamics. The equilibrium concept is inherently analytic: a consistency condition between optimality and aggregation, mirroring the consistency between micro-dynamics and macro-evolution in Laplace's mechanics.


\subsection{Why Lagrangian and Laplacian Systems Don't Fit Directed Acyclic Graphs}

The modern ``causal inference'' literature, exemplified by the work of Judea Pearl and the difference-in-differences approach, relies fundamentally on the concept of a directed acyclic graph (DAG). In a DAG, variables are nodes and causal relationships are directed edges, with the crucial restriction that there are no cycles: one cannot follow the directed edges and return to the starting node. This structure allows for a clear notion of ``upstream'' and ``downstream'' variables, and supports the identification of causal effects through conditional independence relationships.

The dynamical systems studied by Lagrange, Hamilton, and Laplace---and their modern descendants in economics---are fundamentally incompatible with this framework. The incompatibility is not a minor technical issue but a deep structural one.

\subsubsection{Feedback and Cycles}

Consider Hamilton's equations:
\[
\dot{q}_i = \frac{\partial H}{\partial p_i}, \quad \dot{p}_i = -\frac{\partial H}{\partial q_i}.
\]

Each state variable depends on the other state variables, and through time, depends on its own past values. The position $q_i(t)$ at time $t$ depends on the momentum $p_i(s)$ for $s < t$, which in turn depends on $q_i(s')$ for $s' < s$. The system is inherently cyclic: $q$ affects $p$ affects $q$ affects $p$, ad infinitum.

In economic terms, consider a simple optimal consumption-savings problem. Current consumption depends on current wealth. Current wealth depends on past savings. Past savings depend on past consumption. The state variable (wealth) feeds back on itself through the decision rule. There is no acyclic decomposition of this system.

\subsubsection{Simultaneity and Equilibrium}

In an equilibrium system, all variables are determined jointly. Consider a competitive equilibrium: prices affect quantities demanded and supplied, which in turn determine prices through market clearing. The equilibrium price and quantity are simultaneously determined by the intersection of supply and demand.

More formally, in a rational expectations equilibrium, agents' beliefs about the future affect their current actions, which in turn determine the actual future state, which must be consistent with their beliefs. This creates a fixed-point structure:
\[
\text{Beliefs} \to \text{Actions} \to \text{Outcomes} \to \text{Beliefs}.
\]

This is not a causal chain but a consistency condition. The DAG framework, which requires a temporal or logical priority among variables, cannot represent this simultaneous determination.

\subsubsection{The Entire Path Matters, Not Just Discrete Comparisons}

In the variational formulation, the optimal path $q^*(t)$ is the one that minimizes the action functional:
\[
S[q] = \int_{t_0}^{t_1} L(q(t), \dot{q}(t), t) \, dt.
\]

The optimality condition involves the \emph{entire trajectory}, not just values at discrete points. An agent solving an optimal control problem must choose a path that satisfies the Euler--Lagrange equations at every instant, subject to initial and terminal conditions.

The diff-in-diff approach, by contrast, compares outcomes at two discrete time points (before and after treatment) across two groups (treated and control). It cannot capture the continuous-time dynamics, the role of initial conditions, or the global constraints that define optimal trajectories.

\subsubsection{Forward-Looking Behavior}

In Pontryagin's Maximum Principle, the costate variable $\lambda(t)$ satisfies
\[
\dot{\lambda}(t) = -\frac{\partial H}{\partial x}(x^*(t), u^*(t), \lambda(t)),
\]
with terminal condition $\lambda(T) = \frac{\partial \Phi}{\partial x}(x^*(T))$.

The costate evolves \emph{backward in time}. Optimal decisions at time $t$ depend on the anticipated shadow value of the state variable at all future times $s > t$. In economic terms: rational agents are forward-looking. Current investment depends on expected future returns. Current pricing depends on expected future cash flows.

A DAG, with its directed edges representing temporal or causal precedence, cannot represent a variable that ``flows backward'' in time. The causal inference framework treats time as flowing in one direction and assumes that the past affects the future but not vice versa. This is incompatible with rational expectations.

\subsubsection{Conservation Laws and Constraints}

Hamiltonian systems often possess conserved quantities. For an autonomous system (where $H$ does not depend explicitly on $t$), the Hamiltonian itself is conserved:
\[
\frac{dH}{dt} = \sum_i \left( \frac{\partial H}{\partial q_i} \dot{q}_i + \frac{\partial H}{\partial p_i} \dot{p}_i \right) = 0.
\]

In economics, budget constraints, resource constraints, and no-arbitrage conditions play similar roles. These constraints link variables across time and across agents in ways that cannot be represented as a sequence of directed causal edges.

For example, an intertemporal budget constraint requires that the present discounted value of consumption equals the present discounted value of income plus initial wealth:
\[
\int_0^\infty e^{-rt} c(t) \, dt = W_0 + \int_0^\infty e^{-rt} y(t) \, dt.
\]

This is a global constraint on the entire consumption path. Consumption at time $t$ depends on consumption at all other times through this constraint, creating a fully connected structure, not a sparse directed graph.

\subsubsection{The State Space and Phase Space}

In Lagrangian mechanics, the configuration space has dimension $n$ (the number of generalized coordinates), but the state space has dimension $2n$ (positions and velocities). In Hamiltonian mechanics, we work in the phase space $(q, p)$ of dimension $2n$.

The evolution of the system is a flow on this $2n$-dimensional manifold, governed by Hamilton's equations. The symplectic structure $\sum_i dq_i \wedge dp_i$ is preserved by the flow, imposing strong geometric constraints.

Economic dynamics similarly unfold in high-dimensional state spaces. In a macroeconomic model with capital, labor, bonds, and expectations, the state space includes all these variables plus their rates of change or costate variables. The dimensionality and the geometric structure (e.g., constraints imposed by market clearing and optimization) cannot be reduced to a simple graph with a few nodes and edges.

\subsubsection{Why This Matters}

The DAG framework is appropriate for systems where:
\begin{itemize}
\item variables can be ordered such that each variable depends only on ``upstream'' variables,
\item there is a clear distinction between exogenous and endogenous variables,
\item the system can be decomposed into a sequence of conditional distributions,
\item interventions can be modeled as ``surgery'' on individual edges.
\end{itemize}

These conditions fail spectacularly for the dynamical systems central to modern economic theory. In a rational expectations equilibrium, there are no truly exogenous variables (except perhaps deep structural parameters). Everything is endogenous, jointly determined, and linked through dynamic constraints and forward-looking expectations.

To apply causal inference methods to such systems, practitioners must either:
\begin{enumerate}
\item ignore the dynamics and treat the system as if it were static (losing the essence of the economic problem),
\item assume that the treatment is truly exogenous and does not affect other variables (contradicting equilibrium reasoning),
\item work with reduced-form relationships that may be policy-invariant only under very restrictive conditions (violating the Lucas critique).
\end{enumerate}

The structural approach, by contrast, embraces the full complexity of the system. It models the optimization problems, the constraints, the equilibrium conditions, and the dynamics explicitly. Identification comes not from assuming away the feedback and simultaneity, but from imposing the discipline of economic theory: agents optimize, markets clear, beliefs are rational.

This is why Marschak, Koopmans, and their intellectual descendants insisted on structural modeling. They understood, as Lagrange and Laplace did, that complex dynamic systems cannot be reduced to simple causal chains. The proper objects of study are the structural equations, the optimality conditions, and the laws of motion---not isolated ``treatment effects'' extracted from local natural experiments.


\subsection{Hurwicz and Sargent on Causality, Structure, and Identifiability}

The incompatibility between structural economic modeling and causal inference was explicitly recognized by Leonid Hurwicz in his 1967 paper ``On the Structural Form of Interdependent Systems'' \cite{Hurwicz1967}. Hurwicz's analysis, extended and formalized in Sargent's work on equilibrium Markov processes, provides a rigorous foundation for understanding why the concept of causality is problematic in interdependent systems and why it depends fundamentally on the class of anticipated modifications to the system.

\subsubsection{The Identification Problem in Interdependent Systems}

Hurwicz considers a configuration of interdependent components described by a system of equations. In the linear case, such a system takes the form
\[
Ax - b = 0,
\]
where $A$ is the coefficient matrix and $b$ is a vector of constants. The pair $(A, b)$ represents what Hurwicz calls the \emph{behavior pattern} of the configuration.

The fundamental difficulty is that observational data---what Hurwicz calls the ``history'' of the system---cannot uniquely determine the behavior pattern. Any alternative behavior pattern $(C, d)$ related to the true pattern by a nonsingular transformation
\[
C = PA \quad \text{and} \quad d = Pb
\]
will generate the same history. In the jargon of econometrics, the system is \emph{not identifiable} from history alone.

To achieve identifiability, one needs what Hurwicz calls \emph{a priori information}---typically, knowledge that certain variables do not appear in certain equations (i.e., that certain coefficients are zero). The amount of a priori information required depends critically on what predictions one wishes to make.

\subsubsection{Causality is Relative to Anticipated Modifications}

Hurwicz's central insight is that the need for identifying the structural form is not absolute but \emph{relative to the class of anticipated modifications} to the system. He introduces a crucial distinction:

\textbf{No structural change anticipated.} If one only wishes to predict assuming that the future behavior pattern will be the same as in the past, then knowledge of the history suffices. There is no need to identify the structural form, and hence no need for a priori information.

\textbf{Structural modifications anticipated.} If one wishes to predict under hypothetical modifications to the system---for example, the effect of changing the behavior of one component while others remain unchanged---then one must identify the structural form. The degree of identification required depends on the richness of the modification class.

Let $W$ denote the domain of anticipated modifications and let $F_W$ denote the class of behavior patterns that yield correct predictions for all modifications in $W$. Hurwicz shows that:
\begin{itemize}
\item When $W$ consists only of the identity (no modifications), $F_W$ contains all behavior patterns equivalent to the true one---complete non-identifiability.
\item When $W$ is rich enough to include independent modifications of each component, $F_W$ may reduce to a single element (up to normalization)---complete identifiability.
\item The required a priori information $I$ must satisfy $F_I \subseteq F_W$: the class of patterns compatible with the a priori information must be contained in the class permitted by the anticipated modifications.
\end{itemize}

\subsubsection{Structure is Not a Property of the Material System}

Hurwicz emphasizes a point of fundamental importance: the concept of structural form is \emph{not a property of the material system under observation}, but rather \emph{a property of the anticipations of those asking for predictions}. Two researchers who differ in the modifications they are willing to consider will differ in what they accept as structural equations.

This has a crucial implication: if two individuals anticipate different classes of modifications, they may disagree about which equations have ``structural'' or ``lawlike'' status, even though they have access to exactly the same observational data and the same a priori theoretical knowledge.

\subsubsection{Causality Requires Structural Form, But Structure Requires Counterfactuals}

Following Herbert Simon's approach to causality, Hurwicz notes that causal statements like ``$x_1$ is directly causally dependent on $x_2$'' (formalized by Simon as: $x_1$ appears in the second equation but $x_2$ does not appear in the first equation) only make sense if the system is in structural form. But structural form, as Hurwicz has shown, is only defined relative to a class of anticipated modifications.

The crucial implication, stated by Hurwicz on page 238--239, is worth quoting:

\begin{quote}
If, for any reason, one were so to narrow down the class of anticipated modifications as to make the structural form undefined, it would no longer be possible to speak meaningfully of the direction of causality in Simon's sense. But again, the knowledge of this direction would be unnecessary for purposes of prediction.
\end{quote}

This creates a fundamental tension: causality presupposes a decomposition into independently modifiable components, but in a fully interdependent equilibrium system, such a decomposition may not exist. If every variable affects every other variable through equilibrium conditions, there are no ``independently modifiable components,'' and hence no well-defined structural form in the sense required for causal statements.

\subsubsection{Configuration of Independently Modifiable Components}

Hurwicz notes that structural form becomes plausible when the system is a \emph{configuration of independently modifiable components}---that is, when we can modify the behavior pattern of one component while holding others fixed. This is precisely the assumption underlying modern causal inference methods: that we can perform ``surgery'' on individual edges of a causal graph, changing one relationship while leaving others intact.

But this assumption fails in general equilibrium systems. In a competitive equilibrium, we cannot change the demand curve while holding the supply curve and the equilibrium price fixed---they are jointly determined. In a rational expectations equilibrium, we cannot change agents' decision rules while holding the stochastic processes they face and their beliefs fixed---these are linked by the requirement that beliefs equal the mathematical expectation conditional on the equilibrium.

This is the deeper meaning of the Lucas critique: policy interventions change the structure in ways that depend on how the intervention affects agents' optimization problems and equilibrium conditions. The reduced form is not invariant to policy changes because the economy is not a configuration of independently modifiable components.

\subsubsection{Sargent's Formalization Through Equilibrium Markov Models}

Sargent's work on equilibrium Markov processes provides a rigorous formalization and extension of Hurwicz's insights \cite{SargentJPE2025}. The key idea is that \emph{statements about causality are assertions that parameters of a statistical model are invariant with respect to a set of possible interventions}. This is precisely Hurwicz's notion that structural form is defined relative to a modification domain $W$.

Sargent shows how to embed multiple simultaneous notions of causality within equilibrium Markov models through the framework of Markov Decision Problems (MDPs). Each MDP is characterized by a partition of the state space into:
\begin{itemize}
\item A \emph{controllable subspace}: variables the agent can manipulate
\item An \emph{uncontrollable subspace}: variables beyond the agent's control
\end{itemize}

This partition embeds a theory of what causes what from the agent's perspective. Variables in the controllable subspace are causes (instruments); variables in the uncontrollable subspace are effects (outcomes) from that agent's point of view.

In an equilibrium Markov model, there are multiple MDPs operating simultaneously:
\begin{enumerate}
\item The \emph{model builder's MDP}: the researcher's theory of which parameters are invariant to which interventions
\item \emph{Individual agents' MDPs}: each agent's partition of state variables into controllable and uncontrollable
\item \emph{Nature's MDP}: exogenous stochastic processes
\end{enumerate}

Each MDP contains its own theory of causes and effects. The equilibrium reconciles these different perspectives: what appears as an uncontrollable state variable from one agent's perspective may be a controllable decision variable from another agent's perspective, and the equilibrium conditions link these together.

\subsubsection{Granger Causality versus Structural Causality}

Sargent's framework clarifies a crucial distinction that connects back to Hurwicz's analysis: the difference between \emph{Granger causality} and \emph{structural causality}.

Granger causality is a statement about conditional independence within a fixed stochastic process. To say that $x$ Granger-causes $y$ means that past values of $x$ help predict future values of $y$ conditional on past values of $y$. This is purely a property of the joint probability distribution.

Structural causality, by contrast, is a statement about invariance to interventions. To say that $x$ structurally causes $y$ means that the functional relationship from $x$ to $y$ remains stable under a specified class of interventions on $x$.

Sargent illustrates this distinction with a compelling example from German hyperinflation. In time series from hyperinflations, money growth Granger-causes inflation: past money growth helps predict future inflation. But the structural causal relationship may run in the opposite direction: governments print money to finance deficits (inflation causes money growth through a structural budget constraint), while the Granger causality simply reflects timing conventions about when variables are measured.

This example illuminates Hurwicz's point that causality depends on the modification domain. For Granger causality, the modification domain is empty---we only care about prediction within the existing process. For structural causality, the modification domain includes interventions like changes in fiscal policy that affect the budget constraint. The two notions of causality can point in opposite directions.

\subsubsection{Why Fisher's Framework Doesn't Apply to Dynamic Economies}

Sargent emphasizes that the statistical framework underlying modern causal inference---R.A.\ Fisher's theory of randomized experiments with fixed regressors---does not apply to dynamic, interdependent economic systems. Fisher's framework assumes:
\begin{itemize}
\item The researcher controls the treatment assignment
\item Potential outcomes are fixed prior to treatment
\item Treatment does not affect other units (no general equilibrium effects)
\item The data-generating process is stable across treatment and control
\end{itemize}

These assumptions fail in economic systems where:
\begin{itemize}
\item Agents are forward-looking and form expectations about future policies
\item Equilibrium conditions link all agents' decisions and beliefs
\item Interventions change the equilibrium, not just one variable holding others fixed
\item The Lucas critique applies: policy changes alter behavioral equations
\end{itemize}

This is why Hurwicz insisted that structure is relative to anticipated modifications. In a dynamic general equilibrium, one cannot perform Fisher-style ``surgery'' on individual relationships while holding everything else constant. The equilibrium conditions ensure that everything depends on everything else.

Sargent's formalization shows that equilibrium Markov models respect this interdependence while still allowing us to make statements about which parameters are invariant to which interventions. The MDPs make explicit what each agent can and cannot control, and the equilibrium conditions determine how these constraints interact.




\subsubsection{The Sims--Sargent Debate as a Lived Illustration of Hurwicz's Principle}

Hurwicz's abstract argument---that structural form, and hence causal language, is defined
only relative to a class of anticipated modifications---found a concrete and unusually
explicit instantiation in a debate about vector autoregressions that unfolded in
macroeconomics during the early 1980s.  The debate is instructive because the two parties
agreed on nearly all the underlying economic theory, yet reached opposite conclusions about
whether economists could give policy advice.  The disagreement was precisely about the
modification domain.\footnote{Lucas (1987) noted that ``Sims (1982) criticizes the rational expectations revolution for destroying or dis-
carding much that was valuable in the name of utopian ideology’” (Lucas 1987, 8, n. 1).
For Lucas’s perspective on tensions between positive and normative economics, read all of Lucas (1987, n. 1).}

In the rational expectations econometrics program, the historical time series are
interpreted as having been generated by a dynamic game in which private agents solved
constrained intertemporal optimization problems while the government followed a more or
less arbitrary stochastic process for its policy variables.  The modification domain is
therefore non-trivial: the researcher entertains the counterfactual of a government that
behaves differently in the future than it did during the sample period.  Within this
modification domain, deep structural parameters---preferences, technologies, information
structures---are invariant across regimes, and it is meaningful to ask which regime is best.
The vector autoregression is interpreted structurally, and causal language, in Hurwicz's
sense, is well-defined relative to these anticipated interventions.

Christopher Sims challenged this asymmetry.  His view, as interpreted in
Sargent (1984), is that the government, no less than private agents, behaved as a
rational expectations intertemporal optimizer during the sample period.  This is a
symmetric and theoretically coherent position: the dynamic game played during the
estimation period is the same one that will be played after it.  But collapsing the
asymmetry has a decisive consequence for the modification domain.  If government behavior
is governed by the same purposes and constraints before and after the sample period, there
is no room to contemplate a change in regime.  The modification domain shrinks to a single
point---the identity, in Hurwicz's notation---and Hurwicz's theorem then implies complete
non-identifiability of structural form.  With no anticipated modifications, there are no
structural parameters to estimate, no causal relationships to assert, and no policy advice
to give.  Vector autoregressions become, as Sargent (1984) puts it, ``tools for
passive, intervention-free forecasting.''

This is not a failure of Sims' position on its own internal terms.  It is, rather, the
Laplacian conclusion applied to macroeconomics: if the universe---including the
government---is already evolving according to optimal, deterministic-in-distribution laws,
then asking what ``causes'' what, or what should be done differently, is as misconceived as
asking Laplace's demon for advice.  The ghost of cause has been fully exorcised, and
normative economics vanishes along with it.

The rational expectations econometrics position avoids this collapse but at the cost of
an internal contradiction that Sargent (1984) acknowledges directly.  If the
researcher's policy recommendations are persuasive and likely to be adopted, then
forward-looking private agents should already have been anticipating such a regime change
during the estimation period.  The original model, which treated government policy as
arbitrary and exogenous, was therefore misspecified.  Conversely, if regime change is not
anticipated, the recommendations lose much of their practical force.  This contradiction is
a live instance of what Hurwicz had formalized abstractly: the structural form is defined
relative to anticipated modifications, but those modifications must themselves be
incorporated into the model of the historical period if agents are truly forward-looking.
The modification domain and the estimation model are not independent objects.

What the Sims--Sargent debate makes vivid, more clearly than any formal example, is that
the question ``can we give causal or normative advice?'' is equivalent to the question
``is there a non-trivial modification domain?''  Sims answered no; the world is already at
an equilibrium of optimizing behavior, and no causal surgery is possible.  Rational
expectations econometrics answered yes, but only by accepting a model asymmetry that
cannot itself be fully rationalized within a world of perfectly forward-looking agents.
Neither resolution is philosophically clean.  Both positions confirm Hurwicz's central
claim: causality is not a property of the data, nor of the material system, but of the class
of interventions the analyst is prepared to entertain.  To name a cause is always,
implicitly, to specify a counterfactual modification.  And in a world of rational
expectations equilibrium, every such counterfactual modifies the equilibrium itself.

\bigskip

\subsubsection{Hurwicz and the Cowles Commission Tradition}

Hurwicz's approach, developed in collaboration with Marschak, Koopmans, and others at the Cowles Commission, and extended through Sargent's formalization, represents a sophisticated understanding of the relationship between economic theory, econometric identification, and prediction. The Cowles Commission researchers understood that:
\begin{enumerate}
\item Identifiability requires a priori information derived from economic theory.
\item Structural form is defined relative to anticipated policy interventions or modifications.
\item In interdependent systems, causal language can be misleading or meaningless.
\item The proper goal is to identify \emph{structural parameters}---those governing agents' preferences, technologies, and constraints---not causal effects.
\item Multiple perspectives on causality coexist in equilibrium models, reconciled through equilibrium conditions.
\end{enumerate}

From this perspective, the modern ``causal inference'' literature represents a step backward. By focusing on treatment effects identified from natural experiments, without embedding these effects in a structural model, researchers lose the ability to predict the consequences of counterfactual interventions different from those observed in the data. The ``credibility revolution'' in applied econometrics purchased credibility for specific empirical estimates at the cost of abandoning the more ambitious goal of understanding the structural mechanisms governing economic behavior.


\subsection{When Can DAG Methods Be Valid? The MDP Perspective}

The preceding critique of DAG-based causal inference might seem to suggest that such methods are never appropriate in economics. But the MDP framework clarified by Sargent's analysis reveals a more nuanced conclusion: DAG methods can be valid within a carefully circumscribed domain, but they break down precisely when we move from individual decision problems to equilibrium systems.

\subsubsection{DAGs Are Valid Within a Single MDP}

From a single decision maker's perspective, the partition of the state space into controllable and uncontrollable subspaces creates a genuine asymmetry that supports causal reasoning with a DAG structure.

Consider an individual agent facing an MDP. At time $t$, the agent observes the state $s_t$, which includes both variables the agent can control (call these $a_t$, the agent's actions) and variables the agent cannot control (call these $x_t$, the exogenous states). The agent chooses $a_t$ optimally given $s_t = (x_t, \text{endogenous variables})$, and the system transitions to state $s_{t+1}$ according to some law of motion that depends on $(s_t, a_t, \varepsilon_t)$ where $\varepsilon_t$ represents stochastic shocks.

This structure has a natural DAG representation:
\[
x_t \to a_t \to s_{t+1}, \quad x_t \to s_{t+1}, \quad \varepsilon_t \to s_{t+1}.
\]

From this agent's viewpoint, interventions are well-defined. We can perform ``surgery'' on the policy rule mapping states to actions---changing the function $a_t = \pi(s_t)$ to some alternative rule $a_t = \pi'(s_t)$---while holding fixed the transition dynamics for the uncontrollable states $x_t$. This is precisely what is done in policy evaluation within dynamic programming: we compare the value functions or outcome distributions under different policy rules, treating the exogenous state process as invariant.

In this single-agent setting, the controllable/uncontrollable partition defines what Pearl calls a ``manipulation'' or ``intervention'': we can change $a_t$ while holding $x_t$ and the transition law fixed. The DAG correctly represents the causal structure from the agent's perspective.

\subsubsection{DAGs Break Down in Equilibrium}

The problem arises when multiple MDPs interact through equilibrium conditions. Suppose there are two agents, A and B, each solving their own MDP:
\begin{itemize}
\item Agent A's state includes variables $(x_A, y_B)$ where $x_A$ represents truly exogenous variables and $y_B$ represents outcomes from Agent B's decisions that A treats as uncontrollable.
\item Agent B's state includes variables $(x_B, y_A)$ where $x_B$ represents truly exogenous variables and $y_A$ represents outcomes from Agent A's decisions that B treats as uncontrollable.
\end{itemize}

From Agent A's perspective, there is a causal flow: $y_B \to a_A \to y_A$. From Agent B's perspective, there is a causal flow: $y_A \to a_B \to y_B$. But equilibrium requires that these be mutually consistent:
\[
y_A = f_A(x_A, y_B, a_A, \varepsilon_A), \quad y_B = f_B(x_B, y_A, a_B, \varepsilon_B).
\]

The overall system now contains feedback loops that violate the acyclicity required for DAGs:
\[
a_A \to y_A \to a_B \to y_B \to a_A \to \cdots
\]

You cannot perform Pearl-style surgery on Agent A's policy rule while holding Agent B's state process fixed, because Agent B's states depend on Agent A's actions through the equilibrium conditions. Any intervention that changes $\pi_A$ will generically change the equilibrium process for $y_B$, which in turn affects the evolution of Agent A's payoff-relevant states.

This is the essence of the Lucas critique in this framework: the ``reduced form'' relationships between observables are not invariant to policy interventions because the equilibrium links everything together. The parameters that are invariant---the structural parameters in Hurwicz's sense---are the primitives of the individual MDPs (preferences, technologies, information structures), not the equilibrium relationships among endogenous variables.

\subsubsection{The Domain of Validity for DAG Methods}

This analysis suggests that DAG-based causal inference methods are valid in the following circumstances:

\textbf{Single-agent problems.} When studying interventions on a single decision maker who faces a genuinely exogenous environment, the MDP framework supports DAG reasoning. Examples include: medical treatments where the patient does not strategically respond; educational interventions in settings where general equilibrium effects on wages and prices can be neglected; individual firm behavior when the firm is small relative to the market.

\textbf{Partial equilibrium.} When the intervention affects a small part of the economy and the feedback effects on the treated units through general equilibrium channels are negligible. For instance, a local minimum wage change in a small jurisdiction might have valid causal interpretation if labor markets are sufficiently segmented and price effects are small.

\textbf{Agent-level behavioral responses.} DAG methods can identify how individual agents respond to changes in their environment, \emph{holding that environment fixed}. This corresponds to estimating components of an agent's policy function $\pi(s)$, treating the process for $s$ as exogenous. Such estimates are building blocks for structural models but do not by themselves answer policy questions that involve equilibrium effects.

\textbf{Exogenous natural experiments.} When nature performs a randomization that is genuinely independent of all endogenous variables in the system, and where the treatment does not trigger equilibrium responses. Such settings are rare in economics but may occur, for example, in natural disasters or weather shocks that affect some regions but not others.

\subsubsection{Where DAG Methods Fail}

Conversely, DAG methods are problematic in the following settings, which comprise much of macroeconomics and large parts of applied microeconomics:

\textbf{General equilibrium effects.} Policy interventions that change prices, wages, interest rates, or other variables that feed back into agents' decision problems. Examples: monetary policy, fiscal policy, trade policy, large-scale social programs.

\textbf{Strategic interactions.} Game-theoretic settings where agents respond to each other's actions. Examples: industrial organization, bargaining, auctions, political economy.

\textbf{Forward-looking behavior.} Environments where agents form expectations about future policy and adjust their current behavior accordingly. The Lucas critique applies with full force: changing policy changes the decision rules, not just outcomes along a fixed decision rule.

\textbf{Dynamic interdependence.} Settings where today's outcomes affect tomorrow's states for multiple interacting agents. Examples: investment and capital accumulation, human capital and education, savings and wealth distribution.

\subsubsection{A Hierarchical Perspective on Causality}

The MDP framework suggests a hierarchical view of causal structure:
\begin{enumerate}
\item \textbf{Micro level (within MDP):} Each agent's MDP has a valid DAG structure from their own perspective. The partition into controllable and uncontrollable defines causal relationships relative to that agent's feasible interventions.

\item \textbf{Equilibrium level (across MDPs):} The system linking multiple MDPs through equilibrium conditions does not have a DAG structure. Feedback loops and simultaneity make the system acyclic only in a dynamic sense (through time), not in a structural sense (within a time period).

\item \textbf{Structural level (deep parameters):} Objects like utility functions, production functions, information structures, and constraints can be invariant across a wide class of interventions, even when behavioral rules and equilibrium relationships are not. These are the proper targets of structural estimation.
\end{enumerate}

This hierarchy explains why DAG methods sometimes appear to ``work'' in applied microeconomics---they may be valid approximations when studying single agents in partial equilibrium---but are more problematic in macroeconomics and policy analysis, where general equilibrium and strategic interactions are central. The issue is not that causal reasoning is impossible in economics, but that it must be embedded within a structural model that respects the equilibrium nature of economic systems. The MDP framework, precisely because it makes explicit what each agent can and cannot control, provides the right foundation for such reasoning.


\subsection{A Personal Reflection on Diff-in-Diff and the Retreat from Structural Economics}

Although they sometimes agreed about important issues involving economic policies, giants of 20th century economics including Milton Friedman, Paul Samuelson, Kenneth Arrow, T.C.\ Koopmans, Jacob Marschak, Gerard Debreu were united in rarely if ever writing about ``causes''. That is because they knew either the state-of-the-art theory or the state-of-the-art econometrics that they had helped create.

Nevertheless, in recent years an econometric literature on ``causal inference''---particularly the proliferation of difference-in-differences (diff-in-diff) designs---has gained prominence in applied economics. In my view, the popularity of this approach in several areas of ``applied economics'' belies the idea of ``progress'' and is an episode of either ``forgetting'' or ``ignoring'' the lessons from the best parts of mid 20th-century economic theory and econometrics. The methodological posture encouraged by diff-in-diff often requires practitioners to turn their backs on, or inadvertently reveal their ignorance of, the lessons taught by the giants who built the foundations of modern economic theory and econometrics in the twentieth century.

The architects of game theory, general equilibrium theory, and the Cowles Commission approach to econometrics---figures such as Koopmans, Marschak, Arrow, and Hurwicz---worked tirelessly to construct a scientific framework in which economic behavior is understood through systems of structural equations, dynamic optimization, and equilibrium reasoning. Their goal was to build models that respect the interdependence of agents, the feedback inherent in strategic interaction, and the stochastic dynamics governing economic environments.

From this perspective, the diff-in-diff paradigm is a retreat from structural thinking. It encourages researchers to search for isolated ``natural experiments'' and to interpret local mean differences as causal effects, often without embedding these effects in a coherent model of behavior, equilibrium, or dynamics. The resulting estimates may be statistically convenient, but they are frequently disconnected from the economic mechanisms that the Cowles Commission tradition sought to illuminate.

Koopmans and Marschak emphasized that econometric estimation must be grounded in a structural understanding of the system being studied. Arrow and Hurwicz taught us that economic outcomes are the result of interactions among optimizing agents operating within institutions and constraints. These insights are incompatible with the view that causal effects can be reliably inferred from simple before-and-after comparisons across groups, without modeling the strategic and equilibrium forces at work.

The current enthusiasm for diff-in-diff econometrics would have dismayed the twentieth-century giants who built the intellectual foundations of our field. They had created a modern economics cast in terms of structural, dynamic, multivariate equilibria. I personally recommend continuing to capitalize on their legacy by continuing to develop and estimate models that reflect the complexity of the systems we study, rather than by embracing methods that ignore it.


\subsection{Unintended Consequences}

The concept of ``unintended consequences'' is among the oldest in social thought: actions undertaken for one purpose produce effects that were neither foreseen nor desired. It has been invoked by legislators, policymakers, and social commentators for centuries, and Milton Friedman (1992) elevated it to a guiding principle for skepticism about government intervention. But what is the formal status of unintended consequences within modern economic theory? The answer depends entirely on which equilibrium concept one adopts, and the contrast illuminates themes that run throughout this paper.

\subsubsection{No Unintended Consequences in a Rational Expectations Equilibrium}

In a rational expectations equilibrium (REE), a ``communism of beliefs'' prevails: all agents---private citizens and government alike---share a common probability model that coincides with the objective data-generating process \cite{Sargent2008}. Formally, if $f(\cdot \mid \theta)$ is the equilibrium density indexed by structural parameters $\theta$, then every agent's subjective density for endogenous variables equals $f(\cdot \mid \theta)$. No agent is surprised by the consequences of any action, whether taken or not, because the common model correctly describes what would happen under every feasible policy.

This communism of beliefs is what makes the powerful theoretical apparatus of intelligent design work: time inconsistency results, reputation models, optimal taxation, and mechanism design all assume that the government and private agents share a correct model of how the economy responds to any contemplated action. When a government solves a Ramsey problem in an REE, it correctly anticipates the consequences of every policy in its feasible set. There are no unintended consequences because the government's model is correct not only along the equilibrium path but also off it---for policies never chosen and experiments never run.

In the Laplacian language of this paper, the REE is the perfect analogue of Laplace's demon applied to economics. Given the structural primitives (preferences, technologies, information, constraints) and the equilibrium concept, the entire future evolution of the system is determined. Every agent is a local Laplacian demon who sees through the full implications of the structural equations. The ghost of ``cause'' has been exorcised because there is nothing more to explain: the structural equations and initial conditions \emph{are} the explanation, and every consequence is, by construction, intended---or at least correctly foreseen.

\subsubsection{Self-Confirming Equilibrium and Incorrect Off-Path Beliefs}

A self-confirming equilibrium (SCE) weakens the communism of beliefs in a precisely targeted way \cite{Sargent2008}. In an SCE, agents' subjective models match the objective data-generating process for events that occur frequently along the equilibrium path, but they may be wrong about events that are rarely or never observed.

Formally, let the true data-generating mechanism be $f(\mathbf{y}, \mathbf{v} \mid \theta)$ and the government's approximating model be $f(\mathbf{y}, \mathbf{v} \mid \hat{\theta})$. An SCE requires the data-matching condition
\[
f\bigl(\mathbf{y}, v(h \mid \hat{\theta}^{\circ}) \mid \hat{\theta}^{\circ}\bigr) = f\bigl(\mathbf{y}, v(h \mid \hat{\theta}^{\circ}) \mid \theta\bigr),
\]
where $v(h \mid \hat{\theta}^{\circ})$ is the equilibrium policy chosen under the SCE beliefs $\hat{\theta}^{\circ}$. The two densities agree for events that occur under the equilibrium policy, but it is entirely possible that
\[
f(\mathbf{y}, \mathbf{v}' \mid \hat{\theta}^{\circ}) \neq f(\mathbf{y}, \mathbf{v}' \mid \theta)
\]
for alternative policies $\mathbf{v}' \neq v(h \mid \hat{\theta}^{\circ})$ that the government has not tried. The government has correct beliefs about what it has experienced but potentially wrong beliefs about what would happen if it did something different.

The gap between SCE and REE is unimportant for a small agent whose actions have negligible effects on the economy. A small agent's incorrect beliefs about off-equilibrium-path events are harmless precisely because the agent is small: even if he experiments, the consequences for the aggregate economy are negligible. But the gap matters enormously for a large agent like a government. When a government solves a Ramsey problem, it contemplates the consequences of policies not taken---off-equilibrium-path experiments---and its incorrect beliefs about those counterfactuals shape its actual policy choice and therefore the equilibrium path itself.

\subsubsection{Where Unintended Consequences Live}

This is where unintended consequences reside. In an SCE, a government that ventures beyond its historical experience---that tries a new policy it has never tried before---may discover that the actual consequences differ from what its model predicted. The model was confirmed by the data generated under the old policy, but it was never tested against the data that the new policy generates. Milton Friedman's ``law of unintended consequences'' is thus a statement about SCE, not REE: it asserts that governments operate with models that are wrong about precisely those counterfactual events that matter most for policy design.

The adaptive learning literature studied in Sargent (2008) provides a microfoundation for SCE. When agents---including government policymakers---learn about their environment using recursive statistical algorithms, the limiting beliefs converge to an SCE, not necessarily to an REE. Convergence to an SCE means that the learning process generates enough data to discipline beliefs about frequently observed events but not enough to correct beliefs about rarely observed counterfactuals. The government's model fits the data it has seen but may be dangerously wrong about the data it would see under alternative policies.

An extended example from Sargent (2008) makes this vivid. Consider a monetary authority that uses a misspecified Phillips curve to set inflation policy. In the SCE, the authority's estimates of the inflation-unemployment trade-off match the data generated under its  equilibrium policy. But the authority misunderstands how its systematic policy choices shift the Phillips curve through their effects on private agents' inflation expectations. If the authority attempts a permanently lower inflation target---a policy it has never sustained before---it discovers consequences (in this case, more favorable than anticipated) that its model did not predict. The self-confirming equilibrium traps the authority at a suboptimal inflation rate  because the data generated along the equilibrium path do not contain enough variation to reveal the true structural relationship. Only when occasional ``escapes''---unlikely but recurrent sequences of shocks---push the system away from the SCE does the authority temporarily discover a version of the natural rate hypothesis and reduce inflation. These escapes are  temporary precisely because the mean dynamics pull beliefs back toward the SCE, producing the observed pattern of recurrent inflations and stabilizations.

\subsubsection{Connecting to the Themes of This Paper}

The REE--SCE contrast  refracts several themes of this paper through a single prism.

\textbf{Laplace and the exorcism of ``cause.''} In an REE, Laplace's program is fully realized: the structural equations and initial conditions explain everything, and no causal narrative is needed. But in an SCE, something like ``cause'' creeps back in---not as metaphysics, but as a gap between what the decision maker's model predicts and what actually happens when a novel action is taken. The unintended consequence is the discrepancy between the predicted and actual effects of an off-equilibrium-path intervention. This discrepancy exists  because the decision maker's Laplacian model of the world, while confirmed by past experience, is incomplete.

\textbf{Hurwicz and the modification domain.} Hurwicz taught us that structural form is defined relative to a class of anticipated modifications. In an REE, the decision maker's model is structural with respect to every conceivable modification---it correctly predicts the consequences of any feasible policy. In an SCE, the model is structural only with respect to modifications that have been tried. The modification domain implicitly assumed by the decision maker extends beyond the modification domain validated by the data. Unintended consequences arise precisely when the decision maker acts as though a wider modification domain were validated than actually is.

\textbf{DAGs and off-path reasoning.} The diff-in-diff and DAG-based causal inference literatures purport to identify ``causal effects'' from local natural experiments. But the effects so identified are, at best, properties of the SCE in which the  experiment occurred. They provide no guarantee about what would happen under a different policy regime---one that changes the equilibrium and thereby changes the data-generating process. The ``unintended consequences'' problem is precisely the problem that these methods cannot address: what happens when the intervention is large enough, or different enough in kind, to move the system off the path along which the local estimates were obtained.

\textbf{The Sims--Sargent debate.} The contrast between REE and SCE also illuminates the Sims--Sargent debate. Sims's position---that the government is an optimal agent whose behavior should not be treated as exogenous---is the position that the world is already in an REE. In that case, there are no unintended consequences and no scope for policy advice. The rational expectations econometrics position---that the government can be modeled as following a suboptimal rule that could be improved---implicitly assumes an SCE: the government's model is confirmed by the data it has seen but wrong about some counterfactuals. The possibility of giving useful policy advice depends on the existence of this SCE--REE gap.

In sum, the concept of ``unintended consequences'' is vacuous within an REE but substantive within an SCE. The difference is not merely terminological. It determines whether policy advice is possible, whether experimentation can be informative, and whether the ``law of unintended consequences'' is a structural feature of the economic world or merely a confession that we are not yet in an REE. For those of us who believe that actual policymakers operate with models that are at best self-confirming---confirmed by the data they have seen but potentially wrong about the data they have not---the law of unintended consequences is a permanent companion, not a transitional inconvenience to be assumed away.


\appendix


\section{Samuelson on Symmetry, Conservation Laws, and Invariance}

The program of bringing methods from mathematical physics into economics was explicitly launched by Paul A.\ Samuelson in his 1947 \emph{Foundations of Economic Analysis} \cite{Samuelson1947}. Samuelson argued that the unifying principle across economics---from consumer theory to general equilibrium to dynamics---is the method of comparative statics and dynamics grounded in optimization principles analogous to those in classical mechanics.

In \emph{Foundations}, Samuelson showed how economic relationships could be expressed through extremum principles similar to the variational formulations of Lagrange. Demand functions, supply functions, and equilibrium conditions all emerge from optimization subject to constraints. This is an economic counterpart to Lagrange's principle of least action.

But Samuelson went further. Much later in his career, in collaboration with his student Kazuo Sato (who spent many years at NYU), Samuelson explored the deeper connections between symmetry, conservation laws, and invariance in dynamic economic models \cite{SamuelsonSato1984}. This work applied ideas from mathematical physics---particularly the relationship between symmetries and conserved quantities formalized by Emmy Noether---to models of economic growth and capital accumulation.

\subsubsection{Symmetry and Conservation in Physics}

In classical mechanics, Noether's theorem establishes a profound connection: every continuous symmetry of the Lagrangian corresponds to a conserved quantity. For example:
\begin{itemize}
\item \textbf{Time translation invariance} (the laws of physics don't change over time) implies conservation of energy.
\item \textbf{Space translation invariance} (the laws are the same everywhere) implies conservation of momentum.
\item \textbf{Rotational invariance} implies conservation of angular momentum.
\end{itemize}

These conservation laws are not additional assumptions; they are \emph{consequences of the symmetry structure} of the physical system. The conserved quantities are the invariants of the theory.

\subsubsection{Samuelson and Sato on Economic Invariants}

Samuelson and Sato recognized that similar reasoning could be applied to dynamic economic models. In models of optimal growth and capital accumulation, certain symmetry properties of preferences and technology imply the existence of conserved quantities or invariant relationships.

For instance, in a Ramsey growth model with specific functional forms (Cobb--Douglas production, logarithmic utility), certain ratios---such as the capital-output ratio along a balanced growth path---remain constant. These are not arbitrary empirical regularities but rather \emph{structural invariants} that follow from the symmetry properties of the underlying primitives (preferences and technology).

The key insight is that identifying the symmetries of an economic system allows one to isolate what is truly invariant. These invariants---the structural parameters and relationships that remain stable across different configurations or policies---are the proper objects of economic analysis. They correspond to what Hurwicz would later call the ``structural form defined relative to a modification domain.''

\subsubsection{Invariance as the Goal of Structural Modeling}

Samuelson's work on symmetry and conservation laws reinforces the central theme of structural economics: the goal is not to estimate reduced-form relationships that may shift with policy interventions, but rather to identify the deep parameters and invariants that characterize the economic system.

Just as Newton's laws are invariant across different initial conditions and configurations, the structural parameters of an economic model---utility functions, production functions, discount rates, information structures---should be invariant across the class of policy interventions or environmental changes we wish to analyze.

This perspective anticipates the Lucas critique and foreshadows the modern emphasis on structural estimation in macroeconomics. Samuelson understood, from his study of classical mechanics, that the power of a scientific theory lies in identifying what remains invariant under transformations. This is the essence of structural modeling in economics.




\section{Newton, d'Alembert, Lagrange, and Hamilton}

\textbf{Newton: Forces as causes.} Newtonian mechanics is often interpreted as causal because forces appear to ``produce'' accelerations. Newton himself sometimes used causal language, though he also refused to speculate about the cause of gravity. The structure of his laws invites a causal reading: forces act on bodies and thereby change their motion.

Laplace's analytic reformulation removes this causal reading by treating forces as terms in differential equations rather than as metaphysical entities.

\textbf{d'Alembert: Dynamics as equilibrium.} D'Alembert's principle reframes Newton's second law as a condition of equilibrium in an extended configuration space: the sum of applied and inertial forces vanishes for all virtual displacements. This already weakens the causal interpretation of force. Dynamics becomes a matter of balancing terms in an equation, not of one quantity causing another.

\textbf{Lagrange: The elimination of forces.} Lagrange's \emph{M\'ecanique analytique} takes a decisive step away from causal mechanics. By introducing generalized coordinates and the principle of virtual work, Lagrange eliminates explicit forces from the fundamental equations. As he writes in the preface:

\begin{quote}
``On ne trouvera point de figures dans cet ouvrage.''
\end{quote}

or, in English for readers like me,

\begin{quote}
``You will find no figures in this work.''
\end{quote}

The absence of figures symbolizes the elimination of geometric and mechanical intuition in favor of pure analysis. Mechanics becomes a deductive mathematical science. Forces, if they appear at all, are derived quantities, not primitive causes.

\textbf{Hamilton: Mechanics as evolution on phase space.} Hamiltonian mechanics replaces forces with a symplectic flow on phase space. The Hamiltonian function generates time evolution through canonical equations. The fundamental structure is geometric: a flow preserving the symplectic form. Causal language is even less natural here than in Lagrange's formulation. The system evolves according to a geometric law; nothing ``causes'' the evolution beyond the Hamiltonian structure itself.




\section{Mean Field Games}

Mean field games describe strategic interactions among a continuum of agents. Let $u(t, x)$ be the value function and $m(t, x)$ the density of agents.

\subsubsection{Hamilton--Jacobi--Bellman Equation}
\[
-\partial_t u(t, x) = H\left(x, \nabla u(t, x), m(t, \cdot)\right),
\]
with terminal condition $u(T, x) = g(x, m(T))$.

\subsubsection{Fokker--Planck Equation}
\[
\partial_t m(t, x) = -\nabla \cdot \left(m(t, x) D_p H(x, \nabla u(t, x), m(t, \cdot))\right).
\]

\subsubsection{Mean Field Equilibrium}

A pair $(u, m)$ satisfying both PDEs is a mean field equilibrium.


\subsection{Connections to Classical Mechanics}

\subsubsection{Lagrange and MFG}

The HJB equation is a generalized Euler--Lagrange condition for optimal control. Agents choose trajectories that minimize an action functional.

\subsubsection{Hamilton $\to$ MFG}

The HJB equation is the Hamilton--Jacobi equation. The FP equation describes the evolution of a distribution along Hamiltonian characteristics.

\subsubsection{Laplace $\to$ MFG}

The FP equation is a deterministic law of motion for the distribution $m(t, x)$. This is Laplace's worldview applied to populations of optimizing agents.


\section{Hotelling's Pioneering Work on Statistical Manifolds}

In his 1930 abstract ``Spaces of Statistical Parameters'' published in the \emph{Bulletin of the American Mathematical Society} \cite{Hotelling1930}, Harold Hotelling made the groundbreaking observation that parametric families of probability distributions could be viewed geometrically as Riemannian manifolds. This work, appearing in the same decade as the development of differential geometry for general relativity, pioneered the application of non-Euclidean geometry to statistics.

Hotelling proposed treating a parametric statistical model---a family of probability distributions indexed by parameters---as a smooth manifold where each point represents a different probability distribution. He recognized that the Fisher information matrix could serve as a Riemannian metric tensor on this manifold, measuring the ``distance'' between nearby probability distributions in parameter space. Moreover, Hotelling was the first to suggest using the geodesic distance induced by the Fisher metric to measure dissimilarity between statistical models. This Riemannian geodesic distance, now called the Fisher-Rao distance, computes the shortest path between two distributions on the curved manifold.

The significance of Hotelling's insight lies in recognizing that different statistical manifolds have different curvatures. The space of probability distributions is not flat Euclidean space but a curved Riemannian manifold where distances respect the intrinsic geometric structure of the distributions themselves. For instance, while the Fisher-Rao distance between univariate normal distributions can be computed in closed form, there is no closed-form formula for the Fisher-Rao distance between multivariate normal distributions, reflecting the complex curvature of that particular statistical manifold.

Information geometry was thus born independently from econometrician Harold Hotelling (1930) and statistician C.R.\ Rao (1945), both motivated by the mathematical curiosity of considering a parametric family of probability distributions as a Riemannian manifold equipped with the Fisher metric tensor. Hotelling's work laid dormant for over a decade until Rao independently rediscovered similar ideas in his 1945 work on statistical distance. This eventually led to the modern field of information geometry, systematically developed by Chentsov, Amari, and others from the 1960s onward.

The geometric approach to statistics inaugurated by Hotelling has proven remarkably fruitful. By viewing statistical models as manifolds, researchers can apply the powerful machinery of differential geometry---geodesics, curvature, parallel transport---to problems of statistical inference, hypothesis testing, and model comparison. This perspective reveals deep connections between statistical theory and the geometry of curved spaces, connections that continue to generate new insights in both statistics and information theory.

\begin{thebibliography}{9}


\bibitem{Hotelling1930}
Hotelling, H. (1930). Spaces of statistical parameters (abstract). \emph{Bulletin of the American Mathematical Society}, 36, 191.

\bibitem{Hurwicz1967}
Hurwicz, L. (1967). On the structural form of interdependent systems. In E.~Nagel, P.~Suppes, and A.~Tarski (Eds.), \emph{Logic, Methodology and Philosophy of Science: Proceedings of the 1960 International Congress} (pp.~232--239). Stanford University Press.


\bibitem[Sargent(1984)]{Sargent1984}
Sargent, T.~J. (1984).
\newblock Autoregressions, expectations, and advice.
\newblock \emph{American Economic Review}, 74(2):408--415.

\bibitem{Sargent2008}
Sargent, T.~J. (2008). Evolution and intelligent design. \emph{American Economic Review}, 98(1), 5--37.

\bibitem{SargentJPE2025}
Sargent, T.~J. (2025). Macroeconomics after Lucas. \emph{Journal of Political Economy}, 133(11), 3363--3413.

\end{thebibliography}


\end{document}
